% This heading begins below a continued chronology table. Seed the first mark
% with the material at the top of that page, then leave the appendix mark for
% the following page.
\runninghead{Scriptural Date and Location}
\section{Appendix: Scope and Qualifications}
\runninghead{Appendix: Scope and Qualifications}

This study concerns the universal 1962 Roman Missal formulary for the Seventeenth Sunday after Pentecost on 20 September 2026. It includes all ten proper elements and identifies the Credo and Trinity Preface appointments; it does not reproduce the complete Ordinary, expanded prayer conclusions, or Preface. No particular local calendar has been specified. The dated French institutional Ordos corroborate the universal occurrence without supplying an unidentified parish's local overlay.

The Latin is a normalized transcription of the Vatican 1962 typical edition, checked during research against its facsimile and compared with Pustet's 1862 public-domain antecedent. Stress marks and typographical ligatures are normalized; the target's wording and printed extent control. Cummiskey's anonymous 1861 human English supplies the three orations, including its historical phrasing concerning future offenses. The scriptural English is the registered American 1899 Douay--Rheims canonical study text, not a complete English version of the adapted liturgical Latin. Its retained transcription was read; no new collation of the 1899 print is claimed. These inherited texts remain public domain under the recorded source dispositions; the complete modern facsimiles and website presentations are not relicensed by this work.

The biblical contexts examined are Matthew 22, Ephesians 4, Daniel 9, and Psalms 32, 75, 101, and 118, with Matthew 21:23 and 26:1--2 locating the Gospel scene. Psalm citations use the Missal's numbering; common modern equivalents are given with the texts. The Date cells in \textit{Scriptural Date and Location} are generated from \path{research/chronology.toml}, preserving alternatives and unresolved states. The admitted evidence establishes no numerical modern critical date for Daniel or Ephesians and no numerical narrated-event date for the Gospel. A separately attributed comparison gives the NABRE's post-A.D.~70 boundary for the Greek Gospel, without a precise endpoint or any union with the default traditional alternatives. The reception is bounded to the exact loci cited below. Chrysostom and Augustine are principally read in historical English; the NPNF Psalms are abridged. Jerome's Latin and Bellarmine's historical English are retained transcriptions, not newly collated critical texts. Bellarmine's translator omitted philological discussion, variant readings, and original psalm prefaces. Their contributions here are paraphrases. Thomas's Ephesians commentary was researched in the retained 1857 Latin facsimile; the English \textit{Summa} witness includes disclosed electronic-editor modifications.

The three whole-formulary readings are editorial syntheses grounded in these witnesses. No Father is credited with composing or interpreting this precise Sunday as a whole. Thomas's sacramental teaching illuminates the Secret and Postcommunion; no direct ancient exposition of their exact wording or that of the Collect was established in the bounded research. Catena-only attributions and the unverified additional leads listed in the research scope do not support this study. Biblical speakers retain their distinct settings, and personal sin is not attributed to Christ through the voice of his members.

The source research received an independent workflow review before authoring. This document is a study aid, not an official liturgical book, critical edition, or ecclesiastically approved commentary. Text verification, search limits, and exact source dependencies are recorded in \path{propers/verified.md} and \path{research/scope.md}; the interpretation evidence map is \path{research/interpretations.md}. Completed production and artifact reviews are recorded separately in \path{research/production-review.md}.

\section*{References}
\addcontentsline{toc}{section}{References}
\begin{enumerate}
\item \textit{Missale Romanum}, Vatican typical edition (1962), pp.~398--400, nos.~1602--1611; p.~LI, September calendar; p.~XVIII, general rubric 111(b); p.~293, Trinity Preface appointment. \href{https://media.churchmusicassociation.org/pdf/missale62.pdf}{CMAA facsimile}. Antecedent Latin: Pustet, Ratisbon (1862), pp.~341--342, \href{https://archive.org/details/bub_gb_E7sPAAAAIAAJ}{Internet Archive facsimile}, leaves 426--427. The editions have distinct textual-control and rights roles.
\item \textit{The Roman Missal Translated into the English Language for the Use of the Laity} (Philadelphia: Eugene Cummiskey, 1861), pp.~431--433, three orations for the Seventeenth Sunday after Pentecost; \href{https://archive.org/details/romanmissaltran00churgoog}{facsimile}, PDF pp.~440--442. Anonymous human translation; underlying historical text public domain.
\item \textit{Douay--Rheims Bible}, American 1899 edition, registered eBible transcription: Matt.~21:23; 22; 26:1--2; Eph.~1:1; 2:11--18; 4; Dan.~9; Pss.~32,75,101,118. Exact appointed verses and canonical/liturgical distinctions accompany the texts above. Public-domain study English.
\item John Chrysostom, \textit{Homilies on Matthew}, hom.~71, opening exposition of Matt.~22:34--46, trans. George Prevost, revised M.~B. Riddle, NPNF1, vol.~10 (1888), retained CCEL text. \href{https://www.newadvent.org/fathers/200171.htm}{Comparison reading locator}. The later moral continuation is not independently used here.
\item John Chrysostom, \textit{Homilies on Ephesians}, hom.~9, exposition of Eph.~4:1--3; hom.~10, opening exposition of 4:4--5; hom.~11, exposition of common gifts and diverse ministries, NPNF1, vol.~13, retained CCEL text. \href{https://www.newadvent.org/fathers/230109.htm}{Homily 9}, \href{https://www.newadvent.org/fathers/230110.htm}{10}, \href{https://www.newadvent.org/fathers/230111.htm}{11}, comparison locators.
\item Augustine, \textit{Expositions on the Psalms}, NPNF1, vol.~8, retained CCEL English: Ps.~32, first exposition, especially paragraphs 6,12; Ps.~75, paragraphs 3,9--13; Ps.~101, paragraphs 1--3; Ps.~118, paragraphs 34--36,123--124,136. Modern English headings are \href{https://www.newadvent.org/fathers/1801033.htm}{33}, \href{https://www.newadvent.org/fathers/1801076.htm}{76}, \href{https://www.newadvent.org/fathers/1801102.htm}{102}, and \href{https://www.newadvent.org/fathers/1801119.htm}{119}. Numeric paragraphs refer to this abridged English witness, not to Latin sermon numbers.
\item Augustine, \textit{De doctrina christiana}, I.22.20--21 and I.30.31--33, NPNF1, vol.~2, retained CCEL English; \href{https://www.newadvent.org/fathers/12021.htm}{book I comparison locator}. \textit{De consensu evangelistarum}, II.73.141--142, NPNF1, vol.~6, retained CCEL English; \href{https://www.newadvent.org/fathers/1602273.htm}{comparison locator}. The latter witness's final Gospel-name discrepancy is not used to establish a harmonization.
\item Jerome, \textit{Commentaria in Danielem}, II, native comments on Dan.~9:17,18,20,21, Migne, PL~25 (1845), retained Latin transcription, paragraph containing 9:13--21. There is no separate comment on v.~19 in that paragraph. No modern English translation is quoted.
\item Thomas Aquinas, \textit{Super Epistolam ad Ephesios}, ch.~4, lectures 1--2, in \textit{In omnes S. Pauli Apostoli epistolas commentaria}, vol.~2 (Leodii: Dessain, 1857), pp.~312--317. Retained facsimile pp.~316--321. \href{https://isidore.co/aquinas/Eph4.htm}{Modern comparison locator}; its English is not the source of the Latin-based paraphrases here.
\item Thomas Aquinas, \textit{Summa theologiae}, III, q.~79, aa.~1--2,6, corpus, and a.~6 ad 1, English Dominican translation in the retained \href{https://www.gutenberg.org/ebooks/19950}{Gutenberg Tertia electronic edition} (2006). Authorities cited within Thomas are not treated as independently checked witnesses.
\item Robert Bellarmine, \textit{A Commentary on the Book of Psalms}, trans. John O'Sullivan (1866), Ps.~32 ad 6,12; Ps.~75 ad 11--12; Ps.~101 opening commentary; Ps.~118 ad 124. Retained historical-text derivative of the \href{https://www.ecatholic2000.com/bellarmine/commentary-on-psalms.shtml}{eCatholic2000 presentation}; attributed paraphrase only. The opening Ps.~101 verse corresponds to Vulgate 101:2, and Ps.~75 commentary 11--12 to Vulgate 75:12--13.
\item FSSP France, \href{https://www.fssp.fr/ordo-du-mois/}{Ordo du mois}, September 2026, entry for 20 September; ICRSP France, \href{https://icrspfrance.fr/ordo.php}{Ordo}, entry for 20 September 2026. Entries checked during research on 17 September 2026; institutional corroboration of universal occurrence, not a specified local Ordo.
\item Chronology witnesses, retained New Advent article texts in the \textit{Catholic Encyclopedia}: F.~E. Gigot, \href{https://www.newadvent.org/cathen/04621b.htm}{``Book of Daniel''}, vol.~4 (1908), ``Authorship and date of composition''; P.~Ladeuze, \href{https://www.newadvent.org/cathen/05485a.htm}{``Epistle to the Ephesians''}, vol.~5 (1909), ``To whom addressed'', ``Date and place of composition'' and ``Authenticity''; F.~Prat, \href{https://www.newadvent.org/cathen/11567b.htm}{``St.~Paul''}, vol.~11 (1911), ``Chronology'', final table; E.~Jacquier, \href{https://www.newadvent.org/cathen/10057a.htm}{``Gospel of St.~Matthew''}, vol.~10 (1911), ``Destination'' and ``Date and place of composition''; A.~Durand, \href{https://www.newadvent.org/cathen/14530a.htm}{``The New Testament''}, vol.~14 (1912), ``Origin'', paragraph on Matthew's Aramaic original.
\item \textit{New American Bible Revised Edition}, introductions to \href{https://bible.usccb.org/bible/psalms/0}{Psalms} (superscriptions and broad chronological limit) and \href{https://bible.usccb.org/bible/matthew/0}{Matthew} (the separately identified post-A.D.~70 critical boundary and its probabilistic qualification), official USCCB web edition, inspected 21 September 2026. Protected introductions summarized, not reproduced.
\end{enumerate}
